\section{At What Hour You Know Not}

Of the hour of death only one thing is certain, that it is hidden. Christ says it in the same discourse from which the 1962 Missal reads on the Last Sunday after Pentecost: ``Watch ye therefore, because you know not what hour your Lord will come. But this know ye, that, if the goodman of the house knew at what hour the thief would come, he would certainly watch and would not suffer his house to be broken open. Wherefore be you also ready, because at what hour you know not the Son of man will come.''\endnote{Matthew 24:42--44, Douay--Rheims (Challoner).} He ends the parable of the ten virgins with the same command: ``Watch ye therefore, because you know not the day nor the hour.''\endnote{Matthew 25:13, Douay--Rheims (Challoner).} The words speak first of his coming at the end of the world, and the Church has always heard in them also the coming of the Lord to each soul at its death. The thief does not announce himself. A man who will not watch until he hears him at the door will not watch at all.

The Church turns this warning into prayer. Among the devotions for the sick in the \emph{Manual of Prayers} that the Third Plenary Council of Baltimore ordered for the laity, the faithful ask:

\begin{quote}
O God, who hast doomed all men to die, but hast concealed from all the hour of their death: grant that I may pass my days in the practice of holiness and justice, and that I may be made worthy to quit this world in the peace of a good conscience and in the embrace of Thy love. Through Christ our Lord. \emph{Amen.}\endnote{\emph{A Manual of Prayers for the Use of the Catholic Laity}, prepared and published by order of the Third Plenary Council of Baltimore (New York: Catholic Publication Society, 1889), Devotions for the Sick, 469.}
\end{quote}

\noindent The prayer asks for nothing about the hour. It asks to be ready at whatever hour God has kept to himself.

\subsection{The race toward death}

Augustine saw that the road to the hidden hour is already being travelled. ``For no sooner do we begin to live in this dying body, than we begin to move ceaselessly towards death.'' Every day taken from life is taken from the time that remains, ``so that our whole life is nothing but a race towards death, in which no one is allowed to stand still for a little space, or to go somewhat more slowly, but all are driven forwards with an impartial movement, and with equal rapidity. For he whose life is short spends a day no more swiftly than he whose life is longer.''\endnote{Augustine, \emph{De civitate Dei} XIII.10, Dods.} The young man and the old run at the same pace. The only difference is the length of the course, and no runner can see the line.

Benedict wrote the consequence into his Rule. Among the instruments of good works he sets four in a row: \latin{Diem iudicii timere. Gehennam expavescere. Vitam \ae ternam omni concupiscentia spiritali desiderare. Mortem cotidie ante oculos suspectam habere}. To fear the day of judgment, to dread hell, to desire eternal life with all spiritual longing, and to keep death before one's eyes each day as a thing that may come at any time: these are the four last things, in Benedict's own order.\endnote{Rule of Benedict 4.44--47, Latin text (Wikisource transcription); the English is a gloss.} The monk is not asked to imagine death once in a while as a spiritual exercise. He is asked to live as one whose death is already looking at him.

\subsection{The school of the \emph{Imitation}}

The \emph{Imitation of Christ} teaches this lesson with unsparing plainness, and its chapter on the meditation of death deserves to be heard at length. ``Very quickly will there be an end of thee here; take heed therefore how it will be with thee in another world. To-day man is, and to-morrow he will be seen no more. And being removed out of sight, quickly also he is out of mind. \ldots{} Thou oughtest in every deed and thought so to order thyself, as if thou wert to die this day. If thou hadst a good conscience thou wouldst not greatly fear death. It were better for thee to watch against sin, than to fly from death. If to-day thou art not ready, how shalt thou be ready to-morrow?''\endnote{Thomas \`a Kempis, \emph{The Imitation of Christ} I.23.1, trans. William Benham.}

The author does not soften what he has seen. ``If it is a fearful thing to die, it may be perchance a yet more fearful thing to live long. Happy is the man who hath the hour of his death always before his eyes, and daily prepareth himself to die. If thou hast ever seen one die, consider that thou also shalt pass away by the same road.''\endnote{\emph{Imitation} I.23.2, Benham.} And then the counsel that makes the meditation fruitful rather than morbid: ``Happy and wise is he who now striveth to be such in life as he would fain be found in death!'' The things that give ``great confidence of a happy death'' are not feelings at the last hour but habits formed beforehand: contempt of the world, desire of virtue, love of discipline, repentance, obedience, denial of self, and patience under adversity ``for love of Christ.''\endnote{\emph{Imitation} I.23.4, Benham.} The chapter ends where all such meditation ought to end, in prayer: ``To Him direct thy daily prayers with crying and tears, that thy spirit may be found worthy to pass happily after death unto its Lord.''\endnote{\emph{Imitation} I.23.9, Benham.}

\subsection{The meditation of death}

The masters of the following centuries made this remembrance a method. Francis de Sales, writing his \emph{Introduction to the Devout Life} for men and women in the world, gives as its fifth meditation the meditation of death. The soul is to place itself in God's presence, ask his grace, and ``suppose yourself to be on your deathbed, in the last extremity, without the smallest hope of recovery.'' Then it considers. ``One day your soul will quit this body---will it be in summer or winter? in town or country? by day or by night? will it be suddenly or with warning? \ldots{} Alas, of all these things we know absolutely nothing: all that we do know is that die we shall, and for the most part sooner than we expect.''\endnote{Francis de Sales, \emph{Introduction to the Devout Life} I.13, fifth meditation, in the anonymous translation (London: Rivingtons, 1876).} It considers the farewell the soul will take of riches, pleasures, friends, ``husband, wife and child, in short to all creation,'' and last of all of its own body. And it considers the one question that matters: ``when it quits the body the soul must go at once to the right hand or the left. To which will your soul go? what side will it take? none other, be sure, than that to which it had voluntarily drawn while yet in this world.''\endnote{\emph{Introduction} I.13, Rivingtons.} The meditation closes with a resolution and a prayer, and de Sales gives the soul a last instruction for the rest of the day: ``Gather a bouquet of myrrh.''\endnote{\emph{Introduction} I.13, Rivingtons.}

Alphonsus asks for more. In his first consideration he sets the reader beside a body newly dead. ``Imagine that you behold a person who has just expired. Look at that body still laid on the bed, the head fallen on the chest, the hair in disorder and still bathed in the sweat of death, the eyes sunk, the cheeks hollow, the face the colour of ashes.'' Then he follows the mourners home and watches the dead man forgotten, his goods divided, his room given back to the living: ``in the very room in which you shall have breathed forth your soul, and in which you shall be judged by Jesus Christ, others shall dance and eat and play and laugh as before. And where shall your soul then be?''\endnote{Alphonsus Liguori, \emph{Preparation for Death}, Consideration I, first point, Duffy 1843.} The consideration does not end in the grave. It ends at the feet of Christ, in an act of sorrow and a plea for time: ``Lord, I thank you: I accept my death in satisfaction for my sins, and I accept it in the manner in which you shall be pleased to send it. But since you have borne with me till now, wait for me a little longer.''\endnote{\emph{Preparation for Death}, Consideration I, Affections and Prayers, Duffy 1843.}

\subsection{He who lives well}

The whole school can be summed up in Robert Bellarmine's rule. In \emph{The Art of Dying Well}, written in his old age in the quiet of retreat, he set one principle before all the others: ``the general rule, `that he who lives well, will die well,' must be mentioned before all others: for since death is nothing more than the end of life, it is certain that all who live well to the end, die well; nor can he die ill, who hath never lived ill.''\endnote{Robert Bellarmine, \emph{The Art of Dying Well} I.1, trans. John Dalton, digital text (Internet Archive).} Someone will name the good thief, who lived ill and died well. Bellarmine answers that the thief's last hours were a life lived well, short as it was: he defended Christ, rebuked his blaspheming companion, and called upon the Lord before he died. And so ``True, therefore, is the sentence, `He who lives well, dies well;' and, `He who lives ill, dies ill.'\,''\endnote{\emph{The Art of Dying Well} I.1, Dalton.} The hour is hidden so that every hour may be lived as if it were the last. That is the mercy in the concealment. God does not tell us when, in order that we may be ready always. And the readiness cannot be put off past the hour itself, for ``Death is the end of man's earthly pilgrimage, of the time of grace and mercy which God offers him so as to work out his earthly life in keeping with the divine plan, and to decide his ultimate destiny.''\endnote{CCC 1013; compare CCC 1021 and Hebrews 9:27.}
